\subsection{Thermal Expansion and Thermal Stress}

Thermal expansion is the tendency of a solid to change its dimensions when its temperature changes. In most rocks, heating causes expansion while cooling causes contraction, because the average distances between atoms increase with temperature. If a rock specimen were completely unrestrained, this change in size would produce only thermal strain. In reality, however, rock materials are usually constrained by neighboring grains, bedding planes, fractures, joints, or surrounding formations, so free expansion does not occur without resistance. As a result, a temperature change may produce not only deformation, but also stress \cite{boley1960,nowacki1986,jaeger2007}.

A convenient measure of temperature change is the temperature perturbation,
\begin{equation}
\label{eq:theta-def}
\theta = T - T_0,
\end{equation}
where \(T\) is the current temperature and \(T_0\) is a reference temperature, for example the initial ambient temperature of the rock. Thus, \(\theta\) represents the change in temperature from the reference state. Thermal strain and thermal stress depend on temperature differences rather than on absolute temperature alone, so using \(\theta\) or \(\Delta T\) emphasizes the relative change from an initial state \cite{boley1960,nowacki1986}.

Under free expansion, the linear thermal strain \(\varepsilon_{\mathrm{th}}\) produced by a small temperature change is
\begin{equation}
\label{eq:eps-thermal}
\varepsilon_{\mathrm{th}} = \alpha \Delta T,
\end{equation}
where \(\alpha\) is the coefficient of linear thermal expansion and \(\Delta T = T - T_0\) is the temperature change. The thermal strain \(\varepsilon_{\mathrm{th}}\) is dimensionless, while \(\alpha\) has units of inverse temperature, usually \(\mathrm{K^{-1}}\). A positive \(\Delta T\), corresponding to heating, produces positive thermal strain or expansion. A negative \(\Delta T\), corresponding to cooling, produces contraction. This relation describes how the material would deform if its thermal expansion were not resisted by external or internal constraints \cite{boley1960,jaeger2007}.

If a rock is heated but its expansion is fully constrained, thermal strain induces stress. In a simple one-dimensional example with fixed boundaries, the induced thermal stress is
\begin{equation}
\sigma_{\mathrm{th}} = -E\alpha\Delta T,
\end{equation}
where \(\sigma_{\mathrm{th}}\) is the thermal stress and \(E\) is Young's modulus. The negative sign reflects the sign convention in which tensile stress is positive and compressive stress is negative. Under this convention, heating \((\Delta T > 0)\) produces compressive stress because the rock tries to expand but cannot. Conversely, cooling \((\Delta T < 0)\) may produce tensile stress if contraction is restrained. This simple relation shows that thermal stress is proportional to the temperature change and to the stiffness of the material. It should be understood as an illustrative one-dimensional case; real geological media are two-dimensional, heterogeneous, and often only partially constrained \cite{boley1960,jaeger2007}.

In geological media, constraints on thermal expansion are common. At the scale of individual grains, each mineral grain is surrounded by neighboring grains and cement, so its expansion is resisted by contact forces. On large scales, bedding, foliation, or stratification may mechanically limit thermal deformation as adjacent layers can have different stiffness, temperature or thermal expansion behaviour. In situ stresses from overburden and tectonic loading also pre-stress the rock before any additional thermal disturbance occurs. As a result, thermal stress may develop whenever subsurface rocks are heated or cooled under restricted conditions \cite{jaeger2007,turcotte2014}.

Thermal stress is important in several geological and engineering settings. In geothermal reservoirs, injection or circulation of relatively cold fluid in hot rock can alter the local temperature field and change the surrounding stress state. In volcanic and igneous rocks, cooling after emplacement can generate contractional stresses and contribute to cooling joints and columnar fracture patterns. Natural geothermal gradients or engineered heat sources can cause thermal loading in tunnels, boreholes, mines and underground repositories. In fractured rock masses, thermal expansion or contraction may influence crack opening, crack closure, and stress redistribution around discontinuities. These effects are relevant not only for mechanical stability, but also for the conditions under which elastic waves propagate through the rock \cite{robertson1988,turcotte2014,jaeger2007}.

Another source of thermal stress is mineral heterogeneity. Rocks are usually an aggregate of different minerals . Each mineral has its own coefficient of thermal expansion .
Quartz, feldspar, mica, calcite and mafic minerals do not necessarily expand by the same amount for the same change in temperature. Heating (or cooling) the rock may cause one mineral grain to expand (or contract) more than an adjacent grain. This differential expansion (or contraction) creates local stress at grain boundaries, cement contacts, or mineral interfaces. Such mismatches in thermal strain can induce microcracks, change the effective stiffness of the rock. Even within a single anisotropic crystal, thermal expansion may differ between crystallographic directions, adding another source of local stress concentration \cite{robertson1988,mavko2009}.

The linear thermoelastic constitutive relation for an isotropic solid combines mechanical strain and thermal effects. In index notation, the stress tensor \(\sigma_{ij}\) may be written as
\begin{equation}
\sigma_{ij}
=
\lambda \delta_{ij}\varepsilon_{kk}
+
2\mu\varepsilon_{ij}
-
\gamma\delta_{ij}\theta .
\end{equation}
Here, \(\varepsilon_{ij}\) is the small-strain tensor and \(\varepsilon_{kk} = \varepsilon_{11} + \varepsilon_{22}\) is the volumetric strain. The constants \(\lambda\) and \(\mu\) are the Lamé elastic parameters, where \(\mu\) is also the shear modulus. The Kronecker delta \(\delta_{ij}\) equals 1 when \(i=j\) and 0 when \(i \neq j\), so it selects the normal components of the tensor expression. The first two terms,
\(\lambda \delta_{ij}\varepsilon_{kk} + 2\mu\varepsilon_{ij}\), represent the elastic stress caused by mechanical deformation. The last term,
\(-\gamma\delta_{ij}\theta\), represents the stress contribution associated with temperature change, where \(\gamma\) is the thermoelastic coupling coefficient. For an isotropic material, this coefficient is commonly related to the bulk modulus \(K\) and the coefficient of linear thermal expansion \(\alpha\) by
\begin{equation}
\label{eq:gamma-def}
\gamma = 3K\alpha.
\end{equation}
This relation shows that thermal stress depends not only on the magnitude of thermal expansion, but also on the resistance of the material to volumetric deformation \cite{nowacki1986,jaeger2007}.

The same relation can be written in compact tensor form as
\begin{equation}
\sigma
=
\lambda\,\mathrm{tr}(\varepsilon)I
+
2\mu\varepsilon
-
3K\alpha\theta I.
\end{equation}
In this expression, \(\sigma\) is the stress tensor, \(\varepsilon\) is the strain tensor, \(\mathrm{tr}(\varepsilon)\) is the trace of the strain tensor and represents volumetric strain, and \(I\) is the identity tensor. The term \(3K\alpha\theta I\) represents the isotropic thermal part of the stress response. This form makes clear that, in linear isotropic thermoelasticity, stress contains both a mechanical contribution from strain and a thermal contribution from temperature change \cite{boley1960,nowacki1986}.

These classical relations are simplified descriptions of real geological media. Natural rocks may contain pores, fractures, bedding, foliation, mineral fabrics, weathered zones, and fluids, all of which disturb the local distribution of stress and strain. Thermal stress may concentrate near crack tips, grain boundaries, cement contacts, or interfaces between minerals with contrasting expansion behavior. Layered or broken
rock masses may also behave differently depending on the direction of heating, the orientation of discontinuities and the pre-existing stress field. Therefore, the isotropic thermoelastic equations provide a necessary theoretical foundation, but their geological interpretation must account for heterogeneity, anisotropy, and scale \cite{jaeger2007,mavko2009,turcotte2014}.

Thermal expansion and thermal stress are directly relevant to thermoelastic wave propagation. Elastic waves involve time-dependent particle motion and the propagation of stress and strain through a medium. If temperature changes produce additional stress, then the mechanical state through which waves propagate may be altered. Conversely, rocks that are already thermally stressed may have different effective stiffness, crack states, and directional wave behavior compared with unstressed rocks. In heterogeneous geological media, this interaction is controlled by mineral composition, porosity, fractures, thermal conductivity, heat capacity, and elastic moduli. Thermal expansion and thermal stress therefore form a central link between the thermal and mechanical parts of thermoelastic wave propagation in rocks \cite{aki2002,nowacki1986,mavko2009}.